%; whizzy chapter -dvi
% -initex iniptex -latex platex -format platex -bibtex jbibtex -fmt fmt
% 以上 whizzytex を使用する場合の設定。
 
%     Tokyo Debian Meeting resources
%     Copyright (C) 2012 Junichi Uekawa
%     Copyright (C) 2011, 2015 Nobuhiro Iwamatsu

%     This program is free software; you can redistribute it and/or modify
%     it under the terms of the GNU General Public License as published by
%     the Free Software Foundation; either version 2 of the License, or
%     (at your option) any later version.

%     This program is distributed in the hope that it will be useful,
%     but WITHOUT ANY WARRANTY; without even the implied warranty of
%     MERCHANTABILITY or FITNESS FOR A PARTICULAR PURPOSE.  See the
%     GNU General Public License for more details.

%     You should have received a copy of the GNU General Public License
%     along with this program; if not, write to the Free Software
%     Foundation, Inc., 51 Franklin St, Fifth Floor, Boston, MA  02110-1301 USA

%  preview (shell-command (concat "evince " (replace-regexp-in-string "tex$" "pdf"(buffer-file-name)) "&"))

%%ここからヘッダ開始。

\documentclass[mingoth,a4paper]{jsarticle}
\usepackage{monthlyreport}
% 日付を定義する、毎月変わります。
\newcommand{\debmtgyear}{2017}
\newcommand{\debmtgmonth}{5}
\newcommand{\debmtgdate}{20}
% started from zero:
% (let ((year 2013) (month 7)) (+ (* (- year 2005) 12) month -1))
\newcommand{\debmtgnumber}{151}

\begin{document}

\begin{titlepage}
\thispagestyle{empty}
% タイトルページ:編集必要な部分は最初のマクロに飛ばすこと

\vspace*{-2cm}
第\debmtgnumber{}回 東京エリア Debian 勉強会資料\\
\hspace*{-2cm}
\includegraphics{image2012-natsu/dotdeb.pdf}\\
\hfill{}\debmtgyear{}年\debmtgmonth{}月\debmtgdate{}日

% ここはアップデートすること
% 全角文字にしないとフォントのサイズが合わないので注意
\rotatebox{10}{\fontsize{30}{30} {\gt 特集 ：Debbugs}}\\

\vspace*{-2cm}
\hfill{}\includegraphics[height=6cm]{image200502/openlogo-nd.eps}
\end{titlepage}

\newpage

\begin{minipage}[b]{0.2\hsize}
 \definecolor{titleback}{gray}{0.9}
 \colorbox{titleback}{\rotatebox{90}{\fontsize{80}{80} {\gt デビアン勉強会} }}
\end{minipage}
\begin{minipage}[b]{0.8\hsize}
\hrule
\vspace{2mm}
\hrule
\begin{multicols}{2}
\tableofcontents
\end{multicols}
\vspace{2mm}
\hrule
\end{minipage}

\dancersection{最近のDebian関連のミーティング報告}{杉本 典充}

\subsection{第150回東京エリアDebian勉強会}

2017年4月22日(土)に第150回東京エリアDebian勉強会を開催しました。会場は独立行政法人情報処理推進機構(IPA)さんをお借りして行いました。参加者は12名でした。

発表は、yy\_y\_ja\_jpさんによる「Debianパッケージのカスタマイズについて」でした。

packaging-tutorialパッケージを例に、debianパッケージをカスタマイズするときに必要な情報をまとめています。バージョンの命名規則、changelog、パッケージのビルドコマンド、ビルドオプションの変更、quiltを使ってupstreamのソースにパッチを当てる方法、といったパッケージの修正や変更を行うときの手順を発表しました。

その後、勉強会参加者と同じ会場に集まっていたopenSUSEの方々と懇親を深めました。


\dancersection{事前課題}{杉本 典充}

今回の事前課題は以下です:
\begin{enumerate}
  \item Hack Timeは何をしますか。
\end{enumerate}
この課題に対して提出いただいた内容は以下です。
\begin{multicols}{2}
{\small
\begin{prework}{ kenhys }
  \begin{enumerate}
  \item $BH/I=;qNA$r<jD>$7$7$F8x3+$7$F$*$/(B
  \end{enumerate}
\end{prework}

\begin{prework}{ Charles Plessy }
  \begin{enumerate}
  \item r-bioc-* $B%Q%C%1!<%8$N(BQC
  \end{enumerate}
\end{prework}

\begin{prework}{ dictoss }
  \begin{enumerate}
  \item python$B$N(Bsetuptools$B$N3X=,!"(Bpython$B$N(Bdebian$B%Q%C%1!<%8$N:n$jJ}(B
  \end{enumerate}
\end{prework}

\begin{prework}{ yy\_y\_ja\_jp }
  \begin{enumerate}
  \item DDTSS
  \end{enumerate}
\end{prework}

\begin{prework}{ issei }
  \begin{enumerate}
  \item Debconf$B$N%S%G%*4U>^(B
  \end{enumerate}
\end{prework}

\begin{prework}{ Haruki TSURUMOTO(tSU\_RooT) }
  \begin{enumerate}
  \item $BB?J,(BBTS$B$H%a!<%j%9$rFI$`(B
  \end{enumerate}
\end{prework}

\begin{prework}{ ysaito }
  \begin{enumerate}
  \item ($B2sEz$J$7(B)
  \end{enumerate}
\end{prework}

}
\end{multicols}

\dancersection{Debian Trivia Quiz}{杉本 典充}

Debianの昨今の話題についてのQuizです。

今回の出題範囲は\url{debian-devel-announce@lists.debian.org} や \url{debian-news@lists.debian.org}などに投稿された内容からです。

\begin{multicols}{2}
%; whizzy-master ../debianmeetingresume201311.tex
% $B0J>e$N@_Dj$r$7$F$$$k$?$a!"$3$N%U%!%$%k$G(B M-x whizzytex $B$9$k$H!"(Bwhizzytex$B$,MxMQ$G$-$^$9!#(B
%

\santaku
{debian$B%W%m%8%'%/%H$N%Q%C%1!<%84XO"$N%5!<%P$N$&$A!"(BFTP$B%5!<%P$NDs6!$r0lItGQ;_$9$k$3$H$,H/I=$5$l$^$7$?!#(BFTP$B%5!<%P$N$&$A!"$=$N8e$b7QB32TF/$9$kM=Dj$N%5!<%P$O$I$l$G$7$g$&$+!#(B}
{ftp://ftp.debian.org}
{ftp://security.debian.org}
{ftp://ftp.upload.debian.org}
{C}
{$B!V(BShutting down public FTP services$B!W$H$$$&I=Bj$G%"%J%&%s%9$5$l$^$7$?!#(Bapt-line$B$K;XDj$9$k%W%m%H%3%k$ODd;_M=Dj$N(B2017-11-01$B$^$G$K(B"http://"$B$KJQ99$7$^$7$g$&!#$J$*!"3+H/<T8~$1$N(B"ftp://ftp.upload.debian.org"$B!"(B"ftp://security-master.debian.org"$B$O$=$N8e$b7QB32TF/$9$k$H$N$3$H$G$9!#(B \url{https://lists.debian.org/debian-announce/2017/msg00001.html}$B!#(B}

\end{multicols}


% % (query-replace-regexp "<.*?>" "")
% % (query-replace-regexp "^[	 ]\+" "")

%-------------------------------------------------------------------------------
\dancersection{Debbugsとのつきあいかた}{林健太郎(kenhys)}
%-------------------------------------------------------------------------------

\subsection{はじめに}


%
% 冊子にするために、4の倍数にする必要がある。
% そのための調整
%\dancersection{メモ}{}
\mbox{}\newpage


\vspace*{15cm}
\hrule
\vspace{2mm}
\includegraphics[width=2cm]{image200502/openlogo-nd.eps}
\noindent \Large \bf Debian 勉強会資料\\
\noindent \normalfont \debmtgyear{}年\debmtgmonth{}月\debmtgdate{}日 \hspace{5mm}  初版第1刷発行\\
\noindent \normalfont 東京エリア Debian 勉強会 （編集・印刷・発行）\\
\hrule

\end{document}
